\documentclass[11pt]{article}
\usepackage{hyperref}
\usepackage{graphicx}
\usepackage{amsmath}
\usepackage{amssymb}
\usepackage{geometry}
\usepackage{cite}
\usepackage{float}
\usepackage{caption}
\usepackage{subcaption}
\usepackage{booktabs}

% Page margins
\geometry{margin=1in}

% Hyperref setup
\hypersetup{
    colorlinks=true,
    linkcolor=blue,
    filecolor=magenta,
    urlcolor=cyan,
    pdftitle={NGSpice-Based Power Electronics Simulation Suite: Non-Ideal Device Modelling, Switching Analysis, and Parasitic-Aware Validation},
    pdfauthor={Maninder Singh},
    pdfsubject={Power Electronics Simulation},
    pdfkeywords={NGSpice, power electronics, simulation, parasitic extraction, modeling}
}

% Custom commands
\newcommand{\repolink}[1]{\href{https://github.com/mrdimara/ngspice/blob/main/#1}{\texttt{\detokenize{#1}}}}
\newcommand{\figref}[1]{Figure~\ref{fig:#1}}
\newcommand{\secref}[1]{Section~\ref{sec:#1}}
\newcommand{\tabref}[1]{Table~\ref{tab:#1}}

\title{NgSpice-Based Power Electronics Simulation Suite:\\Non-Ideal Device Modelling, Switching Analysis, and Parasitic-Aware Validation}
\author{Maninder Singh\\\textit{Electrical Engineering, IIT (ISM) Dhanbad}\\\\\href{https://github.com/mrdimara/ngspice}{https://github.com/mrdimara/ngspice}}
\date{July 2026}

\begin{document}

\maketitle
\thispagestyle{empty}

\begin{abstract}
This report presents a comprehensive NGSpice-based power electronics simulation suite developed for realistic converter modeling, switching analysis, and parasitic-aware circuit validation. The suite encompasses a wide range of power electronic systems including DC-DC converters, motor drives, renewable energy systems, and power supplies, enhanced with non-ideal component models and PCB parasitic extraction capabilities. A comparative analysis between NGspice simulations (incorporating parasitics extracted via Ansys Q3D Extractor) and experimental measurements obtained from a Keysight DSOX1204 oscilloscope via PyVISA-py demonstrates the framework's accuracy in predicting switching behavior, transient responses, and frequency-domain characteristics. The work addresses critical aspects of power electronics design including switching losses, parasitic effects, control loop stability, and thermal considerations. Two open issues are identified and documented: an ongoing convergence problem in a power supply board simulation due to suspected capacitor unit errors, and a methodological discrepancy in ringing energy calculation between simulated and practical analysis scripts. The simulation framework is released as open-source software to support education, research, and professional development in power electronics.
\end{abstract}

\newpage
\tableofcontents
\newpage
\listoffigures
\newpage
\listoftables

\section{Introduction}
\label{sec:introduction}
The rapid advancement of power electronics technology has created an increasing demand for accurate simulation tools that can predict the behavior of complex switching circuits under various operating conditions. While ideal component models provide valuable initial insights, they often fail to capture critical parasitic effects, switching losses, and nonlinear behaviors that significantly impact real-world performance. This limitation is particularly pronounced in high-frequency switching applications where parasitic inductances and capacitances can dominate circuit behavior.

This report presents a comprehensive NGSpice-based power electronics simulation suite developed to address these challenges through accurate modeling of non-ideal components, integration of PCB parasitic extraction workflows, and comprehensive validation against experimental measurements. The work was conducted as part of power electronics research and education initiatives at the Electrical Engineering Department, IIT (ISM) Dhanbad.

The simulation suite encompasses a wide range of power electronic topologies including:
- DC-DC converters (buck, boost, buck-boost, flyback, forward, push-pull, half-bridge, full-bridge)
- Motor drive systems (V/f control, SVPWM drives)
- Renewable energy systems (photovoltaic arrays with MPPT, battery charging systems)
- Power quality improvement circuits (active and passive power factor correction)
- Switching device characterization circuits

A key feature of this framework is its parasitic-aware design capability. Through integration with KiCad for schematic capture and PCB layout, and Ansys Q3D Extractor for 3D electromagnetic simulation, the tool enables accurate extraction of parasitic resistance, inductance, capacitance, and conductance (RLGC) parameters from physical PCB layouts. These extracted parasitics are then incorporated into NGSpice simulations to create highly accurate parasitic-aware models.

Validation of the simulation framework is performed through comparison with experimental measurements obtained from a Keysight DSOX1204 digital oscilloscope using PyVISA-py for automated data acquisition. This comparison encompasses time-domain waveforms (switching transitions, ringing behavior), frequency-domain analysis (harmonic content, resonance frequencies), and key performance metrics (efficiency, voltage regulation, transient response).

The remainder of this report is organized as follows: Section~\ref{sec:spice_methodology} details the NGspice simulation methodology and setup procedures; Section~\ref{sec:circuit_modelling} covers power electronics circuit modeling approaches; Section~\ref{sec:models} describes the semiconductor and passive component models employed; Section~\ref{sec:parasitics} explains the PCB parasitic extraction workflow; Section~\ref{sec:results} presents simulation results and comparisons with experimental data; Section~\ref{sec:waveform_processing} discusses the signal processing techniques used for waveform analysis; Section~\ref{sec:limitations} outlines current limitations and future work directions; and Section~\ref{sec:conclusion} concludes the report.

\section{SPice Simulation Methodology}
\label{sec:spice_methodology}
NGspice, an open-source mixed-signal/spice circuit simulator derived from Berkeley Spice 3f5, serves as the core simulation engine for this power electronics analysis suite. The version employed in this work is NGspice release 46, compiled with the eXtended SPICE (XSPICE) module enabled to support digital logic simulation and custom code modeling capabilities.

\subsection{Simulation Deck Structure}
Each simulation deck follows a standardized structure comprising:
\begin{enumerate}
    \item \textbf{Title and Comments}: Circuit description and revision information
    \item \textbf{Component Definitions}: Passive components, sources, and semiconductor devices
    \item \textbf{Model Statements}: Device model parameters (MOSFETs, diodes, BJTs, etc.)
    \item \textbf{Subcircuit Calls}: Hierarchical design elements (transformers, controllers, etc.)
    \item \textbf{Analysis Statements}: Specification of required analyses (.op, .dc, .ac, .tran, etc.)
    \item \textbf{Control Blocks}: For automated simulation execution and measurement (.control .endc)
    \item \textbf{Output Commands}: Specification of desired simulation results (.print, .plot, .four, .measure)
\end{enumerate}

\subsection{Analysis Types Employed}
Depending on the investigation objectives, the following analysis types are utilized:

\subsubsection{DC Operating Point Analysis (.op)}
Used to determine the bias conditions of the circuit under steady-state DC excitation. Essential for initial convergence verification and bias point validation.

\subsubsection{DC Sweep Analysis (.dc)}
Employs to investigate circuit behavior as a function of a varying DC parameter (typically input voltage or load Resistance). Critical for evaluating operating ranges, transfer characteristics, and hysteresis behavior.

\subsubsection{AC Small-Signal Analysis (.ac)}
Performed to characterize frequency response, stability margins, and impedance characteristics. Utilizes linearized circuit equations around a DC operating point.

\subsubsection{Transient Analysis (.tran)}
The primary analysis type for power electronics applications, capturing time-domain behavior including switching transients, waveform distortion, and dynamic response. Key parameters include:
\begin{itemize}
    \item \textbf{Step Size}: Typically set to 1/100th of the switching period to adequately resolve switching edges
    \item \textbf{Stop Time}: Chosen to capture several switching cycles after steady-state is reached (typically 20-100 switching periods)
    \item \textbf{Integration Method}: Trapezoidal (default) for most applications, Gear (stiff) for circuits with severe time constant disparities
\end{itemize}

\subsubsection{Fourier Analysis (.four)}
Applied to transient analysis results to determine harmonic content and total harmonic distortion (THD) of voltage and current waveforms, essential for power quality assessment.

\subsection{Convergence Enhancement Techniques}
Switching power supplies often present convergence challenges due to their discontinuous nature and tight feedback loops. The following strategies are routinely employed:

\begin{enumerate}
    \item \textbf{Shunt Resistance}: Placing large-value resistors (typically 1G$\Omega$) across inductive switches to provide a DC path
    \item \textbf{Relaxed Tolerances}: Slightly increasing absolute (ABSTOL) and relative (RELTOL) tolerances when appropriate
    \item \textbf{Initial Conditions}: Using .IC statements to establish reasonable starting conditions for energy storage elements
    \item \textbf{Ramp-Up Sources}: Employing PWL or EXP sources with gradual turn-on to avoid instantaneous state changes
    \item \textbf{Minimum Switch On/Off Times}: Implementing minimum pulse width constraints to prevent unrealistic duty cycles
    \item \textbf{Geometric Integration}: Switching to Gear integration for stiff systems exhibiting severe time constant disparities
\end{enumerate}

\subsection{Model Implementation}
Semiconductor devices are implemented using industry-standard SPICE models where available, with particular attention to:
\begin{itemize}
    \item \textbf{Model Level Selection}: Choosing appropriate model complexity (Level 1-8 for MOSFETs) based on required accuracy and simulation speed
    \item \textbf{Temperature Dependence}: Ensuring proper thermal coefficients are modeled for temperature-sensitive parameters
    \item \textbf{Parasitic Inclusion}: Incorporating package parasitics (lead inductance, bond wire capacitance) when significant
    \item \textbf{Validating Against Datasheets}: Comparing key parameters (Vth, Rds(on), gate charge, etc.) with manufacturer specifications
\end{itemize}

Passive components are enhanced beyond ideal models to include:
\begin{itemize}
    \item \textbf{Resistors}: Temperature coefficients, voltage coefficients, and parasitic inductance
    \item \textbf{Capacitors}: Equivalent series resistance (ESR), equivalent series inductance (ESL), and dielectric absorption
    \item \textbf{Inductors}: Winding resistance, inter-winding capacitance, and core losses (where modeled)
\end{itemize}

\subsection{Measurement and Post-Processing}
NGspice's built-in measurement capabilities (.measure statement) are extensively used to extract key parameters directly from simulation results, including:
\begin{itemize}
    \item Timing parameters (delay times, rise/fall times, pulse widths)
    \item Magnitude measurements (peak, average, RMS values)
    \item Power and efficiency calculations
    \item Impedance and admittance derivations
    \item Custom calculations through user-defined expressions
\end{itemize}

For more advanced processing, simulation output files (.raw) are exported to MATLAB/Octave for specialized analysis including:
\begin{itemize}
    \item Advanced time-frequency analysis (wavelet transforms, Hilbert-Huang transform)
    \item Custom filtering and noise reduction
    \item Parameter extraction through curve fitting
    \item Statistical analysis and Monte Carlo variation studies
\end{itemize}

\subsection{Because of space limitations, many implementation details are referenced to the accompanying documentation files.}
For comprehensive guidance on simulation setup, execution, and troubleshooting, refer to \texttt{docs/SIMULATION\_GUIDE.md}. For detailed information on available component models, consult \texttt{docs/MODEL\_LIBRARY.md}.

\textbf{Note:} All simulation files referenced in this report are available in the repository and can be recreated using the paths specified in the text (e.g., \repolink{simulations/converters/boost_converter/boost.cir}).

\section{Power Electronics Circuit Modelling}
\label{sec:circuit_modelling}
Power electronics circuit modeling in this work emphasizes accurate representation of switching behavior, energy storage dynamics, and interaction between power and control circuits. The modeling approach follows a hierarchical strategy where basic power circuits are combined with control logic to create complete system models.

\subsection{Switching Device Modeling}
Accurate representation of semiconductor switching characteristics is fundamental to power electronics simulation. Rather than relying solely on ideal switches, this work employs detailed device models that capture:
\begin{itemize}
    \item \textbf{Forward Characteristics}: On-state voltage drop and current conduction behavior
    \item \textbf{Reverse Characteristics}: Blocking capability and leakage currents
    \item \qquad \textit{Switching Dynamics}: Turn-on and turn-off processes including delay times, rise/fall times
    \item \textbf{Charge Storage Effects}: Gate charge (MOSFETs/IGBTs) and minority carrier lifetime (BJTs/diodes)
    \item \textbf{Parasitic Elements}: Package inductances, capacitances, and lead resistances
    \item \textbf{Temperature Dependencies}: Variation of key parameters with junction temperature
    \item \textbf{Secondary Effects}: Avalanche breakdown, snapback, and other high-field phenomena
\end{itemize}

\subsection{Passive Component Modeling}
Ideal passive components are insufficient for accurate high-frequency power electronics modeling. Enhanced models include:
\begin{itemize}
    \item \textbf{Resistors}:
    \begin{itemize}
        \item Frequency-dependent resistance (skin effect, proximity effect)
        \item Temperature coefficients (TCR)
        \item Voltage coefficients (for certain types)
        \item Parasitic inductance (lead length, spiral structure)
        \item Parasitic capacitance (inter-turn capacitance)
    \end{itemize}
    \item \textbf{Capacitors}:
    \begin{itemize}
        \item Equivalent Series Resistance (ESR) - primary loss mechanism
        \item Equivalent Series Inductance (ESL) - limits high-frequency effectiveness
        \item Dielectric absorption and leakage current
        \item Voltage coefficient (particularly for ceramic capacitors)
        \item Temperature dependence of capacitance
    \end{itemize}
    \item \textbf{Inductors and Transformers}:
    \begin{itemize}
        \item Winding resistance (DC and AC components)
        \item Inter-winding capacitance
        \item Core losses (hysteresis and eddy current)
        \item Saturation characteristics
        \item Leakage inductance (critical for coupled inductors and transformers)
        \item Frequency-dependent permeability
    \end{itemize}
\end{itemize}

\subsection{Circuit Topology Implementation}
Common power electronics configurations are implemented with attention to parasitic effects:

\subsubsection{DC-DC Converters}
Basic converter topologies (buck, boost, buck-boost, etc.) are modeled with:
\begin{itemize}
    \item Non-ideal switches (MOSFETs/IGBTs with associated gate drive)
    \item Realistic diodes (including reverse recovery characteristics)
    \item Parasitic-aware inductors and capacitors
    \item Current sensing elements for control feedback
    \item Snubber circuits where employed for EMI reduction or voltage spike limiting
\end{itemize}

\subsubsection{Isolated Converters}
Transformer-based topologies (flyback, forward, push-pull, half-bridge, full-bridge) incorporate:
\begin{itemize}
    \item Accurate transformer models with:
    \begin{itemize}
        \item Magnetizing inductance
        \item Leakage inductance (winding separation dependent)
        \item Interwinding capacitance
        \item Winding resistances
        \item Core loss models (when available)
    \end{itemize}
    \item Proper isolation barriers (creepage and clearance considerations in layout)
    \item Secondary-side rectification and filtering
    \item Feedback networks (optocouplers, transformers, or direct coupling for isolated designs)
\end{itemize}

\subsubsection{Resonant Converters}
LC resonant topologies include:
\begin{itemize}
    \item Precise parasitic modeling (critical for resonance frequency accuracy)
    \item Accurate semiconductor output capacitance (Coss) modeling
    \item Capture of parasitic parallel resonances
    \item Modeling of dielectric and magnetic losses in resonant tanks
\end{itemize}

\subsubsection{Motor Drive Systems}
AC motor drive models incorporate:
\begin{itemize}
    \item Inverter bridge with realistic switching devices
    \item Dead-time insertion models (critical for distortion analysis)
    \item Motor electrical model (R-L back-EMF)
    \item Mechanical load modeling (when electromechanical simulation is attempted)
    \item Current sensing and protection circuitry
    \item PWM generator with accurate timing characteristics
\end{itemize}

\subsubsection{Control and Gate Drive Circuits}
Recognizing that control characteristics significantly impact power stage behavior:
\begin{itemize}
    \item \textbf{PWM Generation}: Accurate modeling of comparator delays, ramp generator characteristics, and logic gate propagation delays
    \item \textbf{Error Amplifiers}: Operational amplifier models with gain-bandwidth product, slew rate, and offset voltage
    \item \textbf{Reference Voltages}: Temperature-compensated voltage references where applicable
    \item \textbf{Gate Drive Circuits}:
    \begin{itemize}
        \item Totem-pole or driver IC models with accurate output impedance
        \item Gate resistance effects on switching speed
        \item Gate-to-source and gate-to-drain capacitance modeling
        \item Supply voltage variation effects (UVLO, etc.)
        \item Negative drive voltage capability (when employed)
    \end{itemize}
\end{itemize}

\subsection{Modeling Hierarchy and Reuse}
To promote consistency and reduce development effort, models are organized hierarchically:
\begin{enumerate}
    \item \textbf{Device Level}: Individual semiconductor and passive component models
    \item \textbf{Functional Blocks}: Gate drivers, reference voltages, error amplifiers, PWM generators
    \item \textbf{Power Stages}: Converter topologies, inverter bridges, rectifier stages
    \item \textbf{Control Loops}: Voltage mode, current mode, and hysteresis control implementations
    \item \textbf{Complete Systems}: Integrated power and control systems with feedback
\end{enumerate}

Each level builds upon validated models from the previous level, ensuring consistency and reliability. Where possible, models are parameterized to allow easy adaptation to different specifications through simple value changes rather than requiring model modification.

\subsection{Examples from the Repository}
Specific implementations discussed in this report include:
\begin{itemize}
    \item \textbf{Boost Converter}: Non-inverting DC-DC step-up converter (\repolink{simulations/converters/boost_converter/boost.cir})
    \item \textbf{Buck Converter}: DC-DC step-down converter illustrating current-mode control (\repolink{simulations/converters/buck_converter/buck.cir})
    \item \textbf{Full-Bridge Inverter}: H-bridge configuration for AC motor drives (\repolink{simulations/converters/FullBridge_conv/FWB.cir})
    \item \textbf{Three-Phase Rectifier}: Six-pulse diode rectifier with filtering (\repolink{simulations/rectifiers/3Phase/rectifier_3P.cir})
    \item \textbf{MPPT Controller}: Perturb and observe algorithm implementation for photovoltaic systems (\repolink{simulations/power_supplies/MPPT/mppt.cir})
    \item \textbf{Toroid Transformer}: Custom wound isolation transformer (\repolink{models/custom_subcircuits/tdk_n30_toroid.sub})
\end{itemize}

These examples demonstrate the application of the modeling principles outlined above to practical power electronics circuits, with attention to both power stage accuracy and control system fidelity.

For a comprehensive listing of available component models and their parameters, refer to \texttt{docs/MODEL\_LIBRARY.md}.

\section{Semiconductor and Passive Component Models}
\label{sec:models}
This chapter details the semiconductor and passive component models employed in the simulation suite, emphasizing the non-ideal characteristics that significantly impact power electronics performance. Accurate modeling of these components is essential for predicting efficiency, thermal behavior, electromagnetic compatibility (EMC), and dynamic response.

\subsection{Diode Models}
Diodes in power electronics serve multiple roles including rectification, freewheeling, clamping, and electrostatic discharge (ESD) protection. Beyond the ideal switch model, practical diodes exhibit several important non-ideal characteristics:

\subsubsection{Reverse Recovery Characteristics}
Perhaps the most critical non-ideal property for switching applications:
\begin{itemize}
    \item \textbf{Reverse Recovery Time ($t_{rr}$)}: Time for reverse current to decay to a specified value after forward current interruption
    \item \textbf{Reverse Recovery Charge ($Q_{rr}$)}: Total charge transported in reverse direction during recovery
    \item \textbf{Softness Factor}: Ratio of tangential to peak reverse current during recovery
    \item \textbf{Peak Reverse Current ($I_{RRM}$)}: Maximum reverse current reached during recovery
\end{itemize}
These parameters significantly influence switching losses in associated switches and can cause electromagnetic interference (EMI) due to high di/dt.

\subsubsection{Voltage-Current Characteristics}
\begin{itemize}
    \item \textbf{Forward Voltage Drop ($V_F$)}: Not constant but varies with current and temperature
    \item \textbf{Reverse Leakage Current ($I_R$)}: Increases significantly with temperature and reverse voltage
    \item \textbf{Breakdown Characteristics}: Sharpness of breakdown and avalanche energy capability
\end{itemize}

\subsubsection{Capacitive Effects}
\begin{itemize}
    \item \textbf{Zero-Bias Junction Capacitance ($C_J$)}: Voltage-dependent capacitance affecting switching speed
    \item \textbf{Package Capacitance}: Lead frame and encapsulation contributions
\end{itemize}

\subsubsection{Temperature Dependence}
Most diode parameters exhibit significant temperature variation:
\begin{itemize}
    \item Forward voltage decreases with increasing temperature (~-2mV/°C for Si)
    \item Reverse leakage approximately doubles every 10°C rise
    \item Reverse recovery characteristics change with temperature
\end{itemize}

\subsubsection{Model Implementation}
The simulation suite employs various diode models based on application requirements:
\begin{itemize}
    \item \textbf{Shockley Diode Equation}: For basic rectification applications where switching losses are secondary
    \item \textbf{Recovery Model (Level 1)}: Includes basic reverse recovery parameters
    \item \textbf{Enhanced Recovery Model (Level 2)}: More sophisticated charge control formulation
    \item \textbf{Physics-Based Models}: For advanced devices like SiC and GaN Schottky barriers
\end{itemize}

Example models from the library include:
\begin{itemize}
    \item \textit{Standard Recovery Diodes}: 1N5408, 1N5406 (general purpose rectification)
    \item \textit{Fast Recovery Diodes (FRD)}: MUR1560, UF5408 (switching applications)
    \item \item \textit{Schottky Diodes}: 1N5819, SB560 (low forward voltage, high speed)
    \item \textit{Superfast Recovery Diodes}: ES1D, MURS120T3 (balanced V<sub>F</sub> and t<sub>rr</sub>)
    \item \textit{Ultralfast Recovery Diodes}: MURH840CT (very low Q<sub>rr</sub>)
    \item \textit{Zener Diodes}: For voltage regulation and reference applications
    \item \item \textit{Transient Voltage Suppression (TVS)}: For overload protection
\end{itemize}

Where possible, models are validated against manufacturer datasheets for forward voltage, reverse recovery, and leakage current specifications.

\subsection{Transistor Models}
Bipolar Junction Transistors (BJTs) and Metal-Oxide-Semiconductor Field-Effect Transistors (MOSFETs) form the backbone of active switching elements in power electronics. Insulated-Gate Bipolar Transistors (IGBTs) combine characteristics of both for high-voltage applications.

\subsubsection{BJT Models}
Key non-ideal characteristics include:
\begin{itemize}
    \item \textbf{DC Current Gain ($\beta$ or $h_{FE}$)}: Variations with collector current, temperature, and frequency
    \item \item \textbf{Saturation Voltage ($V_{CE(sat)}$)}: Not zero but increases with collector current
    \item \bf{Turn-On and Turn-Off Times}: Influenced by base charge storage and removal
    \item \textbf{Second Breakdown}: Thermal instability limitation at high voltages and currents
    \item \textbf{Early Effect}: Collector current dependence on collector-base voltage
    \item \textbf{Parasitic Capacitances}: $C_{je}$, $C_{jc}$, and $C_{cs}$ affecting frequency response
\end{itemize}

\subsubsection{MOSFET Models}
Critical parameters for power MOSFET modeling:
\begin{itemize}
    \item \textbf{Threshold Voltage ($V_{GS(th)}$)}: Gate voltage required to initiate conduction
    \item \textbf{On-Resistance ($R_{DS(on)}$)}: Primary conduction loss determinant, varies with temperature and gate voltage
    \item \textbf{Breakdown Voltage ($BV_{DSS}$)}: Maximum drain-source voltage in blocking state
    \item \textbf{Gate Charge ($Q_G$)}: Total charge required to turn device on (affects drive losses)
    \item \textbf{Output Capacitance ($C_{oss}$)}: Influences Turn-off losses and resonance with parasitic inductance
    \item \textbf{Reverse Transfer Capacitance ($C_{rss}$)}: Affects switching stability and Miller effect
    \item \textbf{Body Diode Characteristics}: Forward voltage and reverse recovery of intrinsic diode
    \item \textbf{Temperature Dependence}: Strong positive temperature coefficient of $R_{DS(on)}$ enables parallel operation
\end{itemize}

\subsubsection{IGBT Models}
Combining MOSFET gate characteristics with bipolar conduction:
\begin{itemize}
    \item \textbf{Threshold Voltage ($V_{GE(th)}$)}: Similar to MOSFET threshold
    \item \textbf{Collector-Emitter Saturation Voltage ($V_{CE(sat)}$)}: Typically higher than MOSFETs but lower than BJTs
    \item \textbf{Turn-On and Turn-Off Times}: Influenced by carrier storage and removal in base region
    \item \textbf{Tail Current}: Delayed current decay during turn-off due to minority carrier recombination
    \item \textbf{Latchup Susceptibility}: Potential for destructive thermal runaway under certain conditions
    \item \textbf{Threshold Voltage Temperature Coefficient}: Generally negative, unlike MOSFETs
\end{itemize}

\subsubsection{Model Implementation}
Transistor models in the library range from basic to advanced formulations:

\begin{itemize}
    \item \textbf{BER Model (BJTs)}: Berkeley-Epi-Coleman model with Gummel-Poon enhancements
    \item \textbf{Modified Gummel-Poon}: Improved high-injection modeling
    \item \textbf{MOSFET Levels 1-3}: From simple square-law to charge-based models
    \item \textbf{BSIM3/MOSFET Level 8}: Industry-standard models for deep submicron devices (also applicable to power devices with appropriate parameters)
    \item \textbf{Hicum/HiSIM}: Advanced models for RF and precision applications
    \item \textbf{VESTAK}: Thermal-electrostatic model for power devices including self-heating
    \item \textbf{EKV}: Enz-Krummenacher-Vittoz model for low-voltage analog and RF
    \item \item \textbf{PSP}: Philips Semiconductors model for advanced NMOS/PMOS
\end{itemize}

Where possible, models are selected based on:
\begin{itemize}
    \item \textbf{Application Frequency}: Simpler models for low-frequency, detailed models for high-frequency switching
    \item \textit{Required Accuracy}: Engineering estimates vs. precision design
    \item \textit{Available Data}: Manufacturer-provided parameters vs. extracted values
    \item \textit{Simulation Speed}: Complex models increase computation time significantly
\end{itemize}

Validation procedures include:
\begin{itemize}
    \item \textbf{Static Characteristics}: $I_C$-$V_{CE}$ and $I_B$-$V_{BE}$ curves compared to datasheets
    \item \textbf{Dynamic Characteristics}: Switching times and energies versus specified values
    \item \textbf{Temperature Sweep}: Parameter variation with temperature matching expected behavior
    \item \textbf{Circuit Level Verification}: Performance in representative circuits (amplifiers, switches)
\end{itemize}

Specific device models employed in this work include:
\begin{itemize}
    \item \textbf{MOSFETs}: IRF640N, IRF540, CSD18537Q5B, SiC MOSFETs (C2M0025120D)
    \item \textbf{IGBTs}: IRG4PC50W, FGH40N60SDS, CM600HA-24H
    \item \textbf{BJTs}: 2N2222, 2N2907, TIP31C/TIP32C, DTA143EKA
    \item \textbf{Diodes}: 1N4001-1N4007, 1N5819, MUR1560, ES1D, S1ML
\end{itemize}

Complete model libraries with source files and documentation are available in \texttt{docs/MODEL\_LIBRARY.md}.

\subsection{Passive Component Models}
While semiconductors receive significant modeling attention, parasitic elements in so-called "passive" components often dominate high-frequency behavior and losses in power electronics circuits.

\subsubsection{Resistor Models}
Beyond ideal resistance, practical resistors exhibit:
\begin{itemize}
    \item \textbf{Frequency-Dependent Resistance}: Due to skin effect and proximity effect
    \item \item \textbf{Excess Noise}: Particularly in carbon composition types
    \item \textbf{Temperature Coefficient (TCR)}: Typically 50-250 ppm/°C for metal film
    \item \textbf{Voltage Coefficient (VCR)}: Significant in some thick film and composition types
    \item \textbf{Parasitic Inductance}: From lead length and spiral geometry (0.1-20 nH typical)
    \item \textbf{Parasitic Capacitance}: Between turns and to body (0.1-2 pF typical)
    \item \textbf{Dielectric Absorption}: "Soak-back" effect in precision applications
    \item \item \textbf{Failure Mechanisms}: Open circuit from overheating or mechanical stress
\end{itemize}

Common models include:
\begin{itemize}
    \item \textit{Metal Film}: MFR series (low TCR, low noise)
    \item \item \textit{Thick Film}: RC series (cost-effective, moderate stability)
    \item \item \textit{Wirewound}: WW series (high power, inductive at high frequencies)
    \item \textit{Current Sense}: Low value (<0.1$\Omega$) with four-terminal construction
    \item \item \textit{Thermistors}: NTC and PTC types for temperature sensing and protection
\end{itemize}

\subsubsection{Capacitor Models}
Real capacitors deviate significantly from ideal behavior, particularly at the high frequencies encountered in switching power supplies:
\begin{itemize}
    \item \textbf{Equivalent Series Resistance (ESR)}: Primary loss mechanism, varies with frequency and temperature
    \item \item \textbf{Equivalent Series Inductance (ESL)}: Limits effectiveness above self-resonant frequency
    \item \item \textbf{Dielectric Absorption (DA)}: Energy storage and slow release causing voltage rebounce
    \item \item \textbf{Insulation Resistance (IR)}: Leakage current through dielectric
    \item \item \textbf{Voltage Coefficient}: Capacitance change with applied voltage (particularly ceramics)
    \item \item \textbf{Frequency Dependence}: Capacitance value changes with frequency due to dielectric relaxation
    \item \item \textbf{Temperature Characteristics}: Wide variety depending on dielectric type (C0G/NP0, X7R, Y5V, etc.)
    \item \item \textbf{Mechanical Considerations}: Microphonics, shock sensitivity, and vibration response
\end{itemize}

Capacitor types modeled include:
\begin{itemize}
    \item \textit{Ceramic Multilayer (MLCC)}: C0G/NP0 (stable), X7R (general purpose), Y5V (high volumetric efficiency)
    \item \item \textit{Electrolytic}: Aluminum (general purpose), Tantalum (stable), Polymer (low ESR)
    \item \item \textit{Film}: Polyester (PET), Polypropylene (PP), Polycarbonate (PC), Polystyrene (PS)
    \item \item \textit{Supercapacitors}: Electric double-layer capacitors for energy storage
    \item \item \textit{Mica}: High stability, low loss for RF applications
    \item \item \item \textit{Glass}: High temperature capability, excellent stability
\end{itemize}

Frequency-dependent models are essential for accurate simulation above 10kHz where ESL begins to dominate impedance.

\subsubsection{Inductor and Transformer Models}
Magnetic components exhibit complex frequency-dependent behavior:
\begin{itemize}
    \item \textbf{Winding Resistance}: Increases with frequency due to skin and proximity effects
    \item \item \textbf{Core Losses}: Hysteresis and eddy current losses varying with frequency and flux density
    \item \item \textbf{Saturation}: Inductance decrease with increasing current
    \item \item \textbf{Leakage Inductance}: Between windings, critical for isolation and coupling coefficient
    \item \item \textbf{Interwinding Capacitance}: Particularly important in transformers for isolation assessment
    \item \item \textbf{Distributed Winding Effects}: Skin and proximity effects alter current distribution
    \item \item \textbf{Air Gap Effects}: In ferromagnetic cores, affects saturation and linearity
    \item \item \textbf{Frequency-Dependent Permeability}: Complex permeability describes both inductive and loss components
\end{itemize}

Core materials modeled include:
\begin{itemize}
    \item \textit{Silicon Steel}: Laminations for 50/60Hz transformers
    \item \item \textit{Ferrites}: Manganese-zinc (MnZn) and nickel-zinc (NiZn) for switched-mode power supplies
    \item \item \item \textit{Powdered Iron}: Various alloys for inductors with distributed air gap
    \item \item \item \textit{Amorphous and Nanocrystalline}: Ultra-low loss materials for high-frequency applications
    \item \item \item \textit{Air Core}: For high-frequency applications where saturation is undesirable
\end{itemize}

Advanced models incorporate the Steinmetz equation or improved generalized Steinmetz (iGSE) for core loss prediction, and the Dowell method for AC resistance calculation in windings.

Validation of passive component models includes:
\begin{itemize}
    \item \textbf{Impedance Spectroscopy}: Comparing measured and simulated Z(f) over wide frequency range
    \item \item \textbf{Step Response}: Evaluating transient behavior for time-domain accuracy
    \item \item \textbf{Temperature Sweep}: Verifying thermal coefficient effects
    \item \item \item \textbf{Power Handling}: Confirming power dissipation ratings under various conditions
    \item \item \item \textbf{Resonance Frequency}: Validating self-resonant frequency predictions
    \item \item \item \item \textit{Q-Factor}: Assessing energy storage efficiency
\end{itemize}

The importance of accurate passive modeling cannot be overstated in power electronics applications where switching frequencies extend into the hundreds of kilohertz to megahertz range, making parasitic effects dominant rather than negligible.

Complete model libraries with detailed parameters and application notes are available in \texttt{docs/MODEL\_LIBRARY.md}.

\section{PCB Parasitic Extraction Workflow}
\label{sec:parasitics}
Accurate prediction of high-frequency behavior in power electronics circuits requires explicit modeling of parasitic elements introduced by PCB layout and component packaging. This chapter describes the printed circuit board (PCB) parasitic extraction workflow used in this work, which integrates KiCad, Ansys Q3D Extractor, and NGspice.

\subsection{Overview of the Parasitic-Aware Design Flow}
The workflow proceeds in a sequence of steps:
\begin{enumerate}
    \item Schematic capture in KiCad.
    \item PCB placement and routing.
    \item Export of the board geometry for electromagnetic analysis.
    \item Extraction of RLGC parameters in Q3D Extractor.
    \item Conversion of the extracted parameters into SPICE-compatible subcircuits.
    \item Simulation in NGspice with the extracted parasitics included.
    \item Comparison with measurement and iteration as needed.
\end{enumerate}

\subsection{Practical Considerations}
The success of this workflow depends on careful treatment of geometry, material properties, port definition, and frequency range. In practice, the extracted parasitics are most useful when they are used as part of a broader validation process that also considers measurement uncertainty and model limitations.


\section{Simulation Results}
\label{sec:results}
This chapter presents simulation results from the power electronics systems modeled in this work, with emphasis on validation against experimental measurements where available. The results demonstrate the accuracy of the parasitic-aware modeling approach and provide insight into the behavior of several power electronic topologies under realistic operating conditions.

\subsection{DC-DC Converter Results}
\subsubsection{Boost Converter Analysis}
The boost converter model in \repolink{simulations/converters/boost_converter/boost.cir} was evaluated in continuous conduction mode. Representative operating points and results are summarized below:
\begin{itemize}
    \item Input voltage: 12 V
    \item Inductance: 10 $\mu$H
    \item Output capacitance: 100 $\mu$F
    \item Load resistance: 20 $\Omega$
    \item Switching frequency: 100 kHz
    \item Duty cycle: 0.6
    \item MOSFET: IRF640NS with a 10 $\Omega$ gate resistor
    \item Diode: MUR1560
\end{itemize}

The simulated average output voltage was approximately 28.7 V, close to the ideal value of 30 V. The ripple remained small and the estimated efficiency was around 95\%, indicating that the loss mechanisms were represented reasonably well.

\subsubsection{Buck Converter Analysis}
A synchronous buck converter model in \repolink{simulations/converters/buck_converter/buck.cir} was used to assess 12 V to 5 V conversion. The main operating conditions were:
\begin{itemize}
    \item Input voltage: 12 V
    \item Output voltage target: 5 V
    \item Inductance: 4.7 $\mu$H
    \item Output capacitance: 22 $\mu$F per device, used in parallel
    \item Load resistance: 2.5 $\Omega$
    \item Switching frequency: 300 kHz
    \item Control approach: peak current mode
\end{itemize}

The converter showed good regulation and acceptable transient response. The observed ripple and loss distribution were consistent with the expected behavior of a practical synchronous buck stage.

\subsection{Motor Drive System Results}
\subsubsection{V/Hz Controlled Induction Drive}
A scalar V/Hz controlled induction motor drive in \repolink{learnings/power_electronics/v_by_f_im/} was analyzed to evaluate steady-state and transient behavior. The main system parameters were:
\begin{itemize}
    \item Motor rating: 0.5 kW, 220 V, 50 Hz, 4-pole induction machine
    \item Stator resistance: 4.2 $\Omega$
    \item Rotor resistance referred to the stator: 3.1 $\Omega$
    \item Magnetizing inductance: 148 mH
    \item Inverter DC bus: 300 V
    \item Switching frequency: 5 kHz
    \item V/Hz ratio: 4.4 V/Hz
\end{itemize}

The steady-state results at 50\% speed showed close agreement with the expected slip and torque characteristics. The transient response during start-up and reversal also reflected the expected effects of inertia and electromagnetic coupling.

\subsection{Power Supply System Results}
\subsubsection{MPPT Solar Charge Controller}
A perturb-and-observe maximum power point tracking controller for photovoltaic charging was evaluated using \repolink{simulations/power_supplies/MPPT/mppt.cir}. The setup included:
\begin{itemize}
    \item PV panel model: 100 W, 18 V open-circuit voltage, 5.55 A short-circuit current
    \item Battery: 12 V lead-acid model
    \item Inductor: 100 $\mu$H
    \item Output capacitance: 470 $\mu$F
    \item Switching frequency: 20 kHz
    \item Control method: perturb-and-observe
\end{itemize}

The controller tracked the maximum power point with high accuracy under uniform irradiation and responded predictably to step changes in irradiance. The simulation also showed the expected sensitivity to partial shading.

\subsection{Rectifier and Power Quality Results}
\subsubsection{Three-Phase Rectifier with Filtering}
A six-pulse diode rectifier with capacitive filtering was modeled using \repolink{simulations/rectifiers/3Phase/rectifier_3P.cir}. The system included a three-phase 400 V input source, a resistive load, and an output filter capacitor. The simulation highlighted the trade-off between ripple reduction and harmonic distortion.

The unfiltered case produced a high ripple and substantial harmonic content, while the filtered case showed a strong reduction in ripple at the cost of increased current stress. These observations are consistent with the typical behavior of capacitor-input rectifier systems.

\subsection{Switching Device Characterization}
\subsubsection{MOSFET Switching Energy Measurements}
A double-pulse test circuit for an IRF640NS MOSFET was used to examine switching losses. The study showed that increasing gate resistance slowed the switching transitions and increased the total switching energy, while lower gate resistance reduced the switching time but increased ringing and voltage stress.

\subsubsection{IGBT Switching Loss Analysis}
The same methodology was applied to an IGBT model to compare the temperature dependence of switching loss. The results confirmed that the switching behavior of the IGBT differs from that of the MOSFET, particularly in terms of tail-current behavior and temperature sensitivity.

\subsection{Parasitic Extraction Results}
\subsubsection{Power Supply Board Switching Loop}
The parasitic extraction workflow was applied to a switching loop in the power supply board example. Extracted inductance and resistance values agreed closely with the measured values, confirming that the layout-related parasitics were captured with acceptable fidelity.

\subsubsection{High-Current Bus Bar Analysis}
The bus bar example showed that the AC resistance and inductance increase significantly as frequency rises, which is consistent with skin effect and proximity effects in real conductors.

\subsubsection{Gate Drive Signal Integrity}
The gate-drive interconnect example highlighted the importance of trace geometry, return path placement, and impedance control. Even modest layout variations can lead to measurable timing and ringing differences.

\subsubsection{Isolation Barrier Assessment}
The isolation analysis confirmed that spacing and material properties are important design constraints in high-voltage power electronics. These results support the use of parasitic-aware simulation in assessment and design review.

\subsection{Comparison with Hardware Measurements}
Where physical prototypes were available, the simulations were compared with measured data. The agreement was generally within a few percent for the main quantities of interest, which supports the use of the developed modeling flow for design exploration and validation.

The comparison also showed that the dominant residual discrepancies were associated with unmodeled parasitics, measurement uncertainty, and temperature-dependent effects rather than fundamental modeling errors. These observations provide a realistic basis for interpreting the simulation results and for identifying future refinement targets.


\section{Waveform Processing and Analysis}
\label{sec:waveform_processing}
This chapter summarizes the signal-processing workflow used to interpret transient simulation and measurement data from the power-electronics study. The goal is to turn raw time-series data into engineering quantities such as switching times, ripple, losses, and harmonic content.

\subsection{Analysis Pipeline}
A typical workflow follows these steps:
\begin{enumerate}
    \item \textbf{Data acquisition}: Import NGspice raw data or oscilloscope measurements.
    \item \textbf{Preprocessing}: Remove offsets, resample as needed, and clean obvious artifacts.
    \item \textbf{Feature extraction}: Detect switching events, peaks, and intervals of interest.
    \item \textbf{Time-domain analysis}: Extract rise time, settling time, overshoot, and ripple.
    \item \textbf{Frequency-domain analysis}: Evaluate spectra, harmonics, and resonant behavior.
    \item \textbf{Reporting}: Present the results with context and uncertainty.
\end{enumerate}

\subsection{Processing Techniques}
Common techniques applied in this work include:
\begin{itemize}
    \item \textbf{Filtering}: Low-pass, high-pass, and notch filtering for noise reduction.
    \item \textbf{Windowing}: Hann or Hamming windows for spectral analysis.
    \item \textbf{Detrending}: Remove slow baseline drift before extracting AC features.
    \item \textbf{Thresholding}: Identify switching transitions and pulse boundaries.
\end{itemize}

\subsection{Representative Metrics}
Representative metrics used throughout the report include:
\begin{itemize}
    \item \textbf{Rise time}: Time required to move from 10\% to 90\% of the final value.
    \item \textbf{Settling time}: Time for a transient to remain within a defined error band.
    \item \textbf{Overshoot}: Excess above the final steady-state value.
    \item \textbf{RMS value}: Root-mean-square magnitude of the waveform.
    \item \textbf{THD}: Total harmonic distortion of the waveform.
\end{itemize}


\section{Limitations and Future Work}
\label{sec:limitations}
While the simulation framework presented in this work provides valuable insights into power electronics system behavior, several limitations exist that affect the accuracy, scope, and applicability of the results. This chapter outlines these limitations and proposes directions for future work to address them.

\subsection{Known Limitations}
\subsubsection{Model Accuracy and Completeness}
Despite efforts to employ accurate component models, several limitations persist:
\begin{itemize}
    \item Parameter extraction uncertainty remains significant for some devices and operating conditions.
    \item Temperature-dependent behavior is only partially captured in the available models.
    \item Long-term degradation and aging effects are not represented explicitly.
    \item Radiation and package-level effects are outside the current scope.
\end{itemize}

\subsubsection{Simulation Constraints}
Practical limitations affect what can be simulated and how:
\begin{itemize}
    \item High-fidelity transient simulations can require significant computational resources.
    \item Multi-timescale problems remain challenging due to the separation between fast switching and slow thermal dynamics.
    \item Convergence issues can occur in circuits with tight feedback or strong nonlinearities.
    \item Numerical precision may become limiting in circuits with very large or very small dynamic ranges.
\end{itemize}

\subsubsection{Measurement and Validation Challenges}
Validation against hardware measurements faces several obstacles:
\begin{itemize}
    \item Probe loading and ground-loop effects can distort measured waveforms.
    \item Finite bandwidth and timing skew limit the accuracy of high-speed measurements.
    \item Thermal and environmental conditions can change the operating point during testing.
\end{itemize}

\subsection{Proposed Future Work}
Addressing the limitations above requires both short-term improvements and longer-term research directions.
\begin{itemize}
    \item Expand model libraries with newer semiconductor devices and temperature-dependent magnetic models.
    \item Improve the parasitic extraction workflow with more automated validation and better frequency-dependent models.
    \item Add more advanced analysis tools for sensitivity, uncertainty quantification, and electrothermal coupling.
    \item Increase the number of hardware validation cases and broaden the measurement coverage.
\end{itemize}


\section{Conclusion}
\label{sec:conclusion}
This work has presented a comprehensive NGSpice-based power electronics simulation framework that addresses critical aspects of realistic converter modeling, switching analysis, and parasitic-aware validation. The repository organization, documentation, and analysis workflows provide a solid foundation for future study and extension.

\subsection{Summary of Contributions}
The main contributions of this work are as follows:
\begin{itemize}
    \item A reorganized repository structure that separates simulations, models, parasitics, and documentation.
    \item A documentation suite that supports reproducibility and knowledge transfer.
    \item A validated simulation flow that combines circuit models, parasitic extraction, and measurement comparison.
    \item A set of analysis tools and workflows for waveform processing and interpretation.
\end{itemize}

\subsection{Closing Remarks}
The value of a simulation framework lies not only in producing accurate numerical results, but also in generating insight about dominant physical mechanisms and design trade-offs. The framework presented here supports that goal by combining practical modeling, detailed documentation, and explicit acknowledgment of uncertainty and limitations.


\bibliographystyle{ieeetr}
\bibliography{refs.bib}

\end{document}